\lecture{A posteriori error and adaptive refinement}{Lecture 9 --- Monday, September 14, 2026}{Math 6660 --- \texttt{9th.pdf}}

\begin{roadmap}[Lecture 9]
\roadq{How does a solver decide \emph{where} to refine?}
\begin{roadlist}
\item The \textbf{residual} $R(U)=f+(aU')'$ measures how badly the computed $U$
  fails the strong equation inside each element. It does \emph{not} shrink under
  refinement --- it stays $O(1)$.
\item What shrinks is its weighted contribution: the element indicator
  $E_j^{2}$ multiplies the residual by the element size squared.
\item \textbf{Equidistribution}: split the global error budget evenly over the
  elements and refine anything carrying more than its share. Repeat.
\end{roadlist}
\payoff{This is literally the adaptive mesh refinement loop in a commercial
solver, and $E_j$ is the per-element indicator it colours on the screen.}
\end{roadmap}

% ==================================================================
\subsection*{Error analysis: where we are}

\begin{center}\small
\renewcommand{\arraystretch}{1.2}
\begin{tabular}{@{}ll@{}}
\toprule
\textbf{a priori} & \textcircled{\scriptsize1} energy norm \quad
                    \textcircled{\scriptsize2} $L^{2}$ norm\\
\textbf{a posteriori} & \textcircled{\scriptsize1} energy norm\\
\bottomrule
\end{tabular}
\end{center}

% ==================================================================
\subsection*{A posteriori error, energy norm}

\begin{equation}\label{eq:l9enorm}
  \norm{e}_{E}^{2}=\int_0^1 a\,e'e'\dx ,
\end{equation}
and from the FEM (Lecture 8),
\begin{equation}\label{eq:l9fem}
  \norm{e}_{E}^{2}
   =\int_0^1 f\,\bigl(e-\Pi_h e\bigr)\dx
   -\int_0^1 a\,U'\bigl(e-\Pi_h e\bigr)'\dx .
\end{equation}
But from interpolation,
\begin{equation}\label{eq:l9interp}
  \underbrace{\norm{h^{-1}\bigl(e-\Pi_h e\bigr)}_{a}}_{O(h^{2})}
  \;\le\;C\,\underbrace{\norm{e'}_{a}}_{O(h)} .
\end{equation}

\medskip\noindent\textbf{Note.} Roughly speaking $R(U)=O(1)$.

\begin{intuition}[why the residual does not shrink]
$R(U)$ is the amount by which the computed $U$ fails to satisfy the strong
equation pointwise. It stays $O(1)$ under refinement --- it does not go to zero.
What goes to zero is its \emph{weighted} contribution, because it gets
multiplied by the interpolation error. The estimator is small not because the
residual is small, but because the residual is paired against something small.
\end{intuition}

% ==================================================================
\subsection*{In practice: the adaptive algorithm}

\begin{enumerate}[label=\textcircled{\scriptsize\arabic*},leftmargin=2.2em,itemsep=5pt]

\item On some initial coarse mesh, compute the FEM solution.

\item Calculate on each element the residual
\begin{equation}\label{eq:l9Ej}
  E_{j}^{2}=\bigl(x_{j}-x_{j-1}\bigr)^{2}
            \int_{x_{j-1}}^{x_{j}}\frac{1}{a(x)}\,R^{2}(U)\dx ,
  \qquad\text{where}\qquad
  R(U)=f+\bigl(aU'\bigr)' .
\end{equation}
We want
\begin{equation}\label{eq:l9target}
  E^{2}=\sum_{j=1}^{M+1}E_{j}^{2}\;\le\;\frac{\epsilon^{2}}{c^{2}} ,
\end{equation}
which is achieved if each element satisfies
\begin{equation}\label{eq:l9perelem}
  E_{j}^{2}\;\le\;\frac{\epsilon^{2}}{c^{2}(M+1)} .
\end{equation}

\item Recalculate the solution on the new mesh, then repeat
\textcircled{\scriptsize2} $+$ \textcircled{\scriptsize3}.

\end{enumerate}

\begin{intuition}[equidistribution]
\eqref{eq:l9perelem} splits the global budget $\epsilon^{2}/c^{2}$ evenly over
the $M+1$ elements. An element carrying more than its share gets refined; one
carrying less is left alone. The converged mesh has roughly \emph{equal error
per element} --- which is why adaptive meshes come out fine near a corner or a
load and coarse in the quiet regions, without anyone specifying that by hand.

Note also the factor $(x_j-x_{j-1})^2$ in \eqref{eq:l9Ej}: halving an element
cuts its indicator by four even at fixed residual, so refinement pays.
\end{intuition}

\begin{mechanics}
This is what a commercial solver's adaptive mesh refinement does. $R(U)$ is the
equilibrium residual --- how badly the computed stress field fails to balance
the applied load inside each element --- and $E_j$ is the per-element error
indicator you see coloured on the screen before remeshing.
\end{mechanics}

\todoq{The notes write $\norm{h^{-1}(e-\Pi_h e)}_a \le C\norm{e'}_a$ with
$O(h^2)$ under the left side and $O(h)$ under the right. With the $h^{-1}$
factor in front the left side should be $O(h)$, matching the right; the
$O(h^{2})$ presumably labels $e-\Pi_h e$ alone, before dividing by $h$. Check
which grouping was meant.}
